
\textbf{Lemme 1.5.2.} \uline{Soient} \(A\) \uline{un anneau}, \(B\) \uline{une} \(A\)-\uline{algèbre de présentation finie}, \(M\) \uline{un} \(B\)-\uline{module de présentation finie}. \uline{On suppose} \(B\) \uline{et} \(M\) \uline{plats sur} \(A\), \uline{et} \uline{que pour tout corps résiduel} \(k\) \uline{de} \(A\), \(M\otimes_A k\) \uline{est de dimension projective} \(\leq n\) \uline{sur} \(B\otimes_A k\) (\uline{ce qui est le cas si} \(B=A[t_0,\ldots,t_r]\)). \uline{Sous ces conditions}, \(M\) \uline{admet une résolution de longueur} \(n\) \uline{par des} \(B\)-\uline{modules projectifs de type fini} :
\[
  0\longrightarrow L_n\longrightarrow\ldots\longrightarrow L_0\longrightarrow M\longrightarrow 0.
\]
Ce lemme résulte immédiatement de EGA IV 12.3.2.

\textbf{Proposition 1.6.} \uline{Soient} \(S\) \uline{un schéma quasi-compact}, \(\underline{E}\) \uline{un} \(\underline{O}_S\)\uline{-Module localement libre de rang} \(r+1\), \(\underline{B}\) \uline{son algèbre symétrique}, \uline{avec sa graduation habituelle}, \(\underline{M}\) \uline{un} \(\underline{B}\)\uline{-Module gradué de présentation finie}, \uline{plat sur} \(S\). \uline{Alors il existe une résolution graduée} :
\[
  0\longrightarrow \underline{L}_{r+1}\longrightarrow\ldots\longrightarrow \underline{L}_0\longrightarrow \underline{M}\longrightarrow 0,
\]
\uline{les} \(\underline{L}_i\) \uline{étant des} \(\underline{B}\)\uline{-Modules gradués qui sont localement libres de type fini} ( \uline{en tant que} \(\underline{B}\)\uline{-Modules gradués}).

Montrons d'abord que pour tout \(\underline{B}\)-Module gradué \(\underline{M}\), satisfaisant aux conditions de la proposition, il existe un homomorphisme gradué surjectif~: \(\underline{L}\to \underline{M}\), où \(\underline{L}\) est un \(\underline{B}\)-Module localement libre de type fini.

Comme \(\underline{M}\) est de type fini, il existe un entier \(N\) tel que \(\underline{M}\) soit engendré sur \(\underline{B}\) par ses composantes de degré \(\leq N\). On pose alors :
\[
  \underline{L}=\bigoplus_{i\leq N} \underline{M}_i\otimes_{\underline{O}_S}\underline{B}(-i);
\]
l'homomorphisme de \(\underline{L}\) dans \(\underline{M}\) se définit grâce à (1.2.1), et il est clair qu'il surjectif.

D'après 1.5.2, \(\underline{M}\) admet une résolution (non graduée) par des \(\underline{B}\)-Modules localement libres de type fini sur chaque ouvert affine. Par suite, si 
